\subsection{Operating System Resource Management}

\subsubsection{Cooperative (Non-Preemptive) vs. Preemptive Multitasking Operating System Architectures}

\textbf{Cooperative} scheduling switches tasks only at voluntary suspension points (yield, blocking I/O, explicit await). It is efficient when switching is rare and predictable, but a CPU-bound loop that never yields can starve peers---the same failure mode as a blocked asyncio event loop.

\textbf{Preemptive} scheduling lets the kernel interrupt running code via timer and device interrupts, then choose another ready context. Native processes on desktop and server OSes are preemptive: an infinite loop cannot monopolize the machine forever. The cost is synchronization: shared mutable state can be torn mid-update unless locks or atomics guard critical sequences.

Python straddles both layers: the CPython \emph{process} is preempted by the OS kernel, while coroutines inside one interpreter are usually scheduled cooperatively by an event loop until they \texttt{await}.

\subsubsection{The CPU Kernel Scheduler, Symmetrical Multiprocessing, and Context Switching Latency Overhead}

The kernel scheduler assigns runnable threads to CPU cores. On multicore (SMP) machines, several native threads may execute truly in parallel, but each core still runs one instruction stream at an instant. A context switch saves registers, stack pointer, program counter, and scheduling metadata, then restores another context---work that consumes cycles without advancing user algorithms.

For C programs, this means even tight native loops are periodically interrupted. For CPython, OS-level preemption can occur between any two bytecode instructions, which is why the Global Interpreter Lock serializes execution of Python bytecode in the standard build despite multiple OS threads existing.

\subsubsection{Structural Differentiation: Heavyweight OS Processes (MMU Isolation) vs. Lightweight Native Threads (Shared Memory Spaces)}

A \textbf{process} owns a private virtual address space enforced by the MMU; a \textbf{thread} shares that space with siblings in the same process. Inter-process communication crosses kernel boundaries; inter-thread communication can race on shared memory without copying. Python's multiprocessing spawns new interpreters in separate processes; threading shares one interpreter heap and therefore interacts with the GIL.

\subsubsection{From OS-Level Scheduling to Runtime-Level Scheduling}

Operating systems schedule native instruction streams. Language runtimes schedule additional entities inside one process: greenlets, generators, coroutines, and asyncio tasks. These reuse the cooperative pattern at a higher abstraction: progress continues until an explicit suspension point, then another runtime task runs.

When Python performs asynchronous I/O, the story is not ``thousands of OS threads blocked in \texttt{read()}'' but \textbf{one native thread} driving a loop that registers sockets/files with non-blocking primitives (\texttt{epoll}, \texttt{kqueue}, IOCP). Coroutine frames live on the heap; at \texttt{await}, the frame yields, the event loop services ready handles, and another coroutine resumes. This circumvents per-connection thread overhead and avoids GIL contention between Python bytecode executions in the common single-threaded asyncio deployment, at the price of requiring cooperative yielding for fairness.

\paragraph{Hotel concierge on the grid (synthesis).}
The OS still provides the city grid: virtual memory districts, preemptive street traffic control, and isolated buildings (processes). CPython constructs its luxury hotel inside the heap: objects as furnished rooms, names as luggage tags, reference counting and cyclic GC as concierge staff. Async/event-loop machinery is the hotel's internal bell system---not a second city, but a scheduling layer that multiplexes waiting guests without assigning each guest their own OS thread.
